\subsection{Build-System und Kompilierung}
\label{subsec:build-system-und-kompilierung}

Die Optimierungsstufe je Pipeline-Stufe wird zur Kompilierzeit über
Preprocessor-Makros gewählt, wodurch für jede Konfiguration ein eigenes,
spezialisiertes Binary ohne Laufzeit-Overhead entsteht. Die Stufen
bauen dabei aufeinander auf, statt unabhängige Alternativen
darzustellen: Bei Gaussian ersetzt Shared Memory zunächst die naive
Faltung durch Tiling, während Pinned Memory zusätzlich den
D→H-Transfer beschleunigt, sodass die höchste Stufe stets beides
kombiniert.

\texttt{build\_release.sh} übersetzt jede Übersetzungseinheit einzeln
mit \texttt{nvcc} und reicht die Stufe über die Umgebungsvariablen
\texttt{GAUSSIAN\_OPT}, \texttt{SOBEL\_OPT} und \texttt{NMS\_OPT} als
\texttt{-D}-Flags weiter, zum Beispiel:

\begin{verbatim}
GAUSSIAN_OPT=1 SOBEL_OPT=0 NMS_OPT=1 ./build_release.sh
\end{verbatim}

Ohne Angabe greifen die Standardwerte (\texttt{GAUSSIAN\_OPT=2},
\texttt{SOBEL\_OPT=1}, \texttt{NMS\_OPT=1}), jeweils die am weitesten
optimierte Variante, während \texttt{WARMUP} (Standard: aktiviert)
einen Aufwärm-Kerneldurchlauf vor der Zeitmessung steuert. Da separable
Device-Compilation (\texttt{-rdc=true}) verwendet wird, ist zudem ein
Device-Link-Schritt (\texttt{nvcc~-dlink}) vor dem finalen Host-Link
nötig, bei dem die Ziel-Architektur (\texttt{-arch=sm\_75}) erneut
angegeben werden muss, da \texttt{nvcc} sonst implizit auf eine ältere
Architektur zurückfällt.

\begin{sloppypar}
    Analog erzeugt \texttt{build\_debug.sh} eine Debug-Variante mit den
    Flags \texttt{-G~-O0}, ohne \texttt{--use\_fast\_math}; alle
    berichteten Messungen stammen ausschließlich aus dem Release-Build.
    Das Ausführen der jeweils gebauten Binary übernehmen die separaten
    Skripte \texttt{run\_release.sh} bzw. \texttt{run\_debug.sh} mit fest
    hinterlegtem Ein- und Ausgabepfad.
\end{sloppypar}